% ======================================================================
% SECTION 2: HR 9502 - Cancer Patient Robot-and-Humanoid Choice Act of 2026
% Revision of California AB 2847 (2026) and FDA AI Regulatory
% Decision-Making Draft Guidance 2025.
% ======================================================================

HR 9502 is a \textit{Revision of California AB 2847 (2026) Patient
Rights and Robotic Safety Act and FDA AI Regulatory Decision-Making
Draft Guidance of 2025}, drawing the chain of prior authority from
the California state-level patient-rights baseline up to the FDA
capability-neutrality draft \cite{fda_ai_regulatory_decision_draft_2025}.
The California baseline is documented at
\texttt{new-trial/site/02-legislation-patient-rights/physical\_ai\_patient\_rights.tex}
and at
\texttt{national-platform/new\_trial\_psl/02\_ab2847\_patient\_rights.tex}
\cite{kawchak_2026_19176370, kawchak_2026_19244918}. HR 9502 implements
the seventh dimension of the operational definition of patient
control introduced in this paper: physical-AI procedure execution
choice. The patient selects which surgical robot, humanoid companion,
or rehabilitation exoskeleton performs the patient's care from the
ten robot categories enumerated at
\texttt{patients/patient\_robot\_instructions\_fixed.tex} pages 1
through 10 and at
\texttt{national-platform/patient\_robot/01\_preamble\_robots1to5.tex}
(surgical, cobot, RT positioning, needle placement, social companion)
and \texttt{02\_robots6to10\_closing.tex} (humanoid, RT motion-
tracking, imaging, steerable needle, rehabilitation exoskeleton)
\cite{kawchak_2026_18810541, kawchak_2026_18778219,
target2025robotic_bronchoscopy}.

The operational definition of patient control across six prior
functional dimensions (trial discovery, enrollment comprehension,
location choice, data and biospecimen governance, toxicity warnings,
AI-exclusion contestability) is supplied by
\texttt{patients/paper/Deep-Research-3/part1\_legal\_baseline.md}.
This paper introduces the seventh dimension that HR 9502 implements
\cite{hhs_45cfr46_protection_human_subjects,
fda_informed_consent_guidance_2023, fda_dct_guidance_2024}.

\subsection{Findings and Declarations}

Congress finds and declares the following.

(a) Ten distinct robot categories are clinically deployed in
oncology care today: surgical robots, collaborative cobots, RT
positioning robots, robotic needle-placement systems, social
companion robots, humanoids, RT motion-management robots, imaging
assistant robots, steerable needle robots, and rehabilitation
exoskeletons \cite{kawchak_2026_18810541}.

(b) The Unification Standard Level (USL) framework provides a
four-dimension scoring sheet (capability, safety, autonomy, trust)
that quantifies each robot's readiness for patient choice on a 0
through 100 scale; the USL framework is documented at
\texttt{national-platform/usl\_standard/01\_preamble\_introduction\_framework.tex}
and \texttt{02\_results\_discussion\_closing.tex}
\cite{kawchak_2026_18778219}.

(c) Direct robotics evidence supporting patient-led robot choice
includes the TARGET multicenter robotic-bronchoscopy trial published
in CHEST 2025 (safety and diagnostic outcomes across multiple sites)
and the NEJM Catalyst Cancer CARE Beyond Walls home-based
chemotherapy study published in 2026 (safety, feasibility, and
patient experience metrics) \cite{target2025robotic_bronchoscopy,
dronca2026cancer_care_beyond_walls}.

(d) The California AB 2847 (2026) Patient Rights and Robotic Safety
Act establishes patient rights at the State level only; HR 9502
extends those rights to national coverage through revision rather
than displacement of the State framework \cite{kawchak_2026_19176370}.

(e) The FDA AI Regulatory Decision-Making Draft Guidance of January
2025 articulates the capability-neutrality principle for AI systems
used in regulatory decision-making; HR 9502 incorporates that
principle into its definitions of qualified robot
\cite{fda_ai_regulatory_decision_draft_2025}.

(f) Foundation-model evidence in computational pathology (Virchow,
Nature Medicine 2024) demonstrates that AI systems can reach clinical-
grade performance across pan-cancer detection; the same pattern
generalizes to robot-executed image-guided procedures
\cite{vorontsov2024virchow}.

\subsection{Definitions}

(a) \billterm{Patient robot choice} means the patient's right to
select which surgical robot, humanoid robot, or companion robot
performs the patient's care, exercised after review of the qualified
robot comparison required by Subsection 2.3(a). The reference list of
robot categories is the ten-category enumeration at
\texttt{patients/patient\_robot\_instructions\_fixed.tex} pages 1
through 10 \cite{kawchak_2026_18810541}.

(b) \billterm{Qualified robot} means a Physical AI system that meets
the Unification Standard Level (USL) threshold for the system's
category, as documented at \texttt{national-platform/usl\_standard/}
and at \texttt{unification/usl/paper/usl\_oncology\_trials.tex}
\cite{kawchak_2026_18778219}.

(c) \billterm{Robot category} means one of the ten categories
enumerated in Subsection 2 of the prior California SB 1042 baseline,
repeated at
\texttt{new-trial/site/01-legislation-authorization/physical\_ai\_trial\_authorization.tex}
and at
\texttt{national-platform/new\_trial\_psl/01\_preamble\_and\_sb1042.tex},
with each category carrying its own USL minimum threshold
\cite{kawchak_2026_19176370}.

(d) \billterm{Capability-neutral} means the bill's definitions apply
to any robot architecture (LLM-driven, foundation-model, agentic
system, or specialized hardware) that meets the USL threshold,
following the second guardrail in
\texttt{patients/paper/Deep-Research-3/part2\_ai\_robotics\_legislation.md}
and the FDA AI Regulatory Decision-Making Draft 2025
\cite{fda_ai_regulatory_decision_draft_2025}.

(e) \billterm{USL transparency report} means the four-dimension USL
scoring sheet (capability, safety, autonomy, trust) for a specific
robot instance, published in machine-readable form by the qualified
site that operates the robot \cite{kawchak_2026_18778219}.

\subsection{Patient Robot Choice Rights}

(a) Right to receive a side-by-side comparison of all qualified
robots for the patient's procedure, including USL scores, prior
outcomes, and named instances. Worked examples drawn from the
single-patient journey include da Vinci Xi at USL 87.5 for lobectomy
(\texttt{patient-journey/stage\_05\_surgery.py}), post-operative
cobot for limb stabilization
(\texttt{patient-journey/stage\_06\_recovery.py}), Franka Emika at
USL 88.75 for pembrolizumab infusion
(\texttt{patient-journey/stage\_07\_immunotherapy.py}), and a
companion robot for pediatric or geriatric surveillance
(\texttt{patient-journey/stage\_09\_surveillance.py})
\cite{kawchak_2026_19119939}.

(b) Right to select a specific robot instance (for example RTPOS-01
versus RTPOS-03) following the per-hour
\texttt{patient\_records.md} naming pattern at
\texttt{new-trial/national-24-7-trial/hour-12/patient\_records.md}
\cite{kawchak_2026_19994945}.

(c) Right to substitute a humanoid robot (Tesla Optimus, Boston
Dynamics Atlas, Agility Digit, or other USL-qualified humanoid
evaluated under \texttt{unification/usl/humanoids/}) for routine non-
surgical tasks (transport, vitals collection, medication delivery)
where the patient prefers reduced human contact
\cite{kawchak_2026_18778219}.

(d) Right to refuse a specific robot instance and request a
substitute at no additional cost to the patient.

(e) Right to receive a USL transparency report for any robot before
consenting to the robot's use \cite{kawchak_2026_18778219}.

(f) Right of appeal to the FDA Oncology Center of Excellence DCT
pathway for any robot denial that the patient believes is
capability-equivalent to the requested robot
\cite{fda_oce_advancing_oncology_dct_2024}.

(g) Right to a humanoid-robot deployment for in-home care over the
course of long therapeutic regimens (the 35 pembrolizumab cycles for
PAT-2026-0042 are a worked example), with a patient-side override
authority that may revert any in-home humanoid deployment to human-
nurse delivery within the same calendar day
\cite{kawchak_2026_19119939}.

\subsection{Implementation}

(a) The Secretary of Health and Human Services, acting through the
Food and Drug Administration, the Agency for Healthcare Research and
Quality, and the Office of the National Coordinator for Health
Information Technology, shall require all qualified Physical AI
oncology trial sites to publish a per-robot USL transparency report
within 12 months of enactment.

(b) The bill is a Revision of California AB 2847 (2026) for national
coverage; State-level protections that exceed the federal floor are
preserved by the preemption rule
\cite{kawchak_2026_19176370}.

(c) Implementation references the operational pattern at
\texttt{national-platform/patient\_robot/02\_robots6to10\_closing.tex}
and the safety-workflow integration at
\texttt{sponsor/final\_paper/scripts/safety/}
\cite{kawchak_2026_19244918, kawchak_2026_19396256}.

\subsection{Reporting and Enforcement}

(a) Annual federal report listing the per-state, per-cancer, per-
robot-category selection rate. Domain 5 of the federal Patient
Control Dashboard (Physical-AI Execution) anchors this report:
patient-selected robot percentage by category, patient-initiated
procedural modification rate (tied to HR 9503), patient-sponsor
direct channel uptake (tied to HR 9505), and the American Leadership
Index ranking (tied to HR 9506) \cite{kawchak_2026_18810541,
kawchak_2026_18778219, kawchak_2026_19244918}.

(b) The AI-decision-tracking infrastructure at
\texttt{sponsor/final\_paper/scripts/core\_agents/} (53 agents listed
in
\texttt{sponsor/final\_paper/scripts/core\_agents/\_\_init\_\_.py})
provides the FDA-auditable substrate \cite{kawchak_2026_19396256}.

(c) Civil penalties harmonize with the FDA AI Regulatory Decision-
Making Draft 2025 capability-neutrality enforcement framework
\cite{fda_ai_regulatory_decision_draft_2025}.

\subsection{Discussion and Limitation}

The bill rests on the assumption that USL scoring is comparable
across vendors. The unification/usl/ paper documents that USL scoring
requires multi-site calibration; humanoid robots in particular have
a thinner clinical evidence base than surgical robots. The findings
above acknowledge this limitation and constrain humanoid deployment
to non-surgical routine tasks until additional prospective evidence
accumulates \cite{kawchak_2026_18778219}.

\subsection{Patient-Control Framings}

For the individual patient, a 58-year-old female NSCLC patient
matching the PAT-2026-0042 profile chooses the da Vinci Xi for her
lobectomy after reviewing USL 87.5 versus the USL 84.0 alternative,
and chooses a Tesla Optimus humanoid for home-based vitals collection
during her 35 pembrolizumab cycles, with a daily revert-to-human-
nurse override available at any point during the regimen
\cite{dronca2026cancer_care_beyond_walls, kawchak_2026_19119939}.

For broad adoption, all Physical AI oncology trial sites must
publish a USL transparency report for every deployed robot instance
within 12 months. The bill positions the United States as Number 1
in the world for patient-controlled robot selection, replacing the
prior provider-only choice pattern under traditional CA AB 2847 with
a national patient-led standard governing the ten robot categories
\cite{kawchak_2026_18810541, kawchak_2026_19176370}.
